\documentclass[11pt,a4paper]{article}

% Packages
\usepackage[utf8]{inputenc}
\usepackage[T1]{fontenc}
\usepackage{amsmath,amssymb,amsfonts}
\usepackage{mathtools}
\usepackage{graphicx}
\usepackage{booktabs}
\usepackage{hyperref}
\usepackage{cleveref}
\usepackage[margin=1in]{geometry}
\usepackage{fancyhdr}
\usepackage{titlesec}
\usepackage{enumitem}
\usepackage{bm}

% Custom commands
\newcommand{\vect}[1]{\bm{#1}}
\newcommand{\mat}[1]{\mathbf{#1}}
\newcommand{\quat}[1]{\mathbf{#1}}
\newcommand{\deriv}[2]{\frac{d#1}{d#2}}
\newcommand{\pderiv}[2]{\frac{\partial #1}{\partial #2}}
\newcommand{\norm}[1]{\left\lVert #1 \right\rVert}
\DeclareMathOperator{\sgn}{sgn}
% Simple unit command (replacement for siunitx)
\newcommand{\si}[1]{\,\mathrm{#1}}
\newcommand{\SI}[2]{#1\,\mathrm{#2}}

% Header/Footer
\pagestyle{fancy}
\fancyhf{}
\rhead{Space State Simulation}
\lhead{Core Physics Documentation}
\rfoot{Page \thepage}

% Title
\title{%
    \textbf{Space State Astrodynamics Simulation} \\
    \large Core Physics and Mathematical Foundations
}
\author{%
    Space State Development Team \\
    \texttt{space\_state v1.0}
}
\date{December 2025}

\begin{document}

\maketitle

\begin{abstract}
This document presents the mathematical foundations and physics models implemented in the Space State astrodynamics simulation package. We derive the equations of motion for spacecraft trajectory propagation, attitude dynamics, propulsion systems, and guidance algorithms. The simulation employs high-fidelity n-body gravitational dynamics in the Solar System Barycentric (SSB) reference frame, quaternion-based attitude representation, and Zero-Effort-Miss/Zero-Effort-Velocity (ZEM/ZEV) guidance for finite burn maneuvers. All equations are presented with sufficient detail for implementation verification and validation.
\end{abstract}

\tableofcontents
\newpage

%==============================================================================
\section{Introduction}
%==============================================================================

The Space State simulation provides a comprehensive framework for spacecraft trajectory analysis, mission planning, and guidance system development. The core physics models include:

\begin{itemize}
    \item N-body gravitational dynamics with JPL ephemeris data
    \item Quaternion-based attitude kinematics and dynamics
    \item Reaction wheel and RCS thruster control allocation
    \item Finite burn guidance using ZEM/ZEV algorithms
    \item Orbital mechanics with Keplerian elements
\end{itemize}

This document serves as the authoritative reference for all physics models implemented in the simulation.

%==============================================================================
\section{Reference Frames}
%==============================================================================

\subsection{Solar System Barycentric Frame (SSB)}

The primary reference frame is the Solar System Barycentric (SSB) frame, which is an inertial frame centered at the barycenter of the solar system. Position and velocity vectors are expressed in astronomical units (AU) and AU/day respectively.

The SSB frame uses the J2000 equatorial coordinate system:
\begin{itemize}
    \item \textbf{X-axis}: Pointing toward the vernal equinox at J2000.0
    \item \textbf{Z-axis}: Perpendicular to the mean equator at J2000.0, pointing north
    \item \textbf{Y-axis}: Completing the right-handed system
\end{itemize}

\subsection{Body-Centered Frames}

For local orbital mechanics, we transform to body-centered frames. Given the SSB state $(\vect{r}_\text{SSB}, \vect{v}_\text{SSB})$ and body ephemeris $(\vect{r}_\text{body}, \vect{v}_\text{body})$:

\begin{align}
    \vect{r}_\text{rel} &= \vect{r}_\text{SSB} - \vect{r}_\text{body} \\
    \vect{v}_\text{rel} &= \vect{v}_\text{SSB} - \vect{v}_\text{body}
\end{align}

\subsection{Unit Conversions}

The simulation maintains internal consistency with the following constants:

\begin{table}[h]
\centering
\begin{tabular}{lll}
\toprule
\textbf{Quantity} & \textbf{Value} & \textbf{Units} \\
\midrule
Astronomical Unit & $1.495978707 \times 10^{8}$ & km \\
Day & $86400$ & seconds \\
Gravitational constant $G$ & $6.67430 \times 10^{-11}$ & $\si{m^3.kg^{-1}.s^{-2}}$ \\
Standard gravity $g_0$ & $9.80665$ & $\si{m.s^{-2}}$ \\
\bottomrule
\end{tabular}
\caption{Fundamental constants used in the simulation.}
\label{tab:constants}
\end{table}

%==============================================================================
\section{Gravitational Dynamics}
%==============================================================================

\subsection{N-Body Problem}

The equation of motion for a spacecraft under the gravitational influence of $N$ celestial bodies is:

\begin{equation}
    \ddot{\vect{r}} = -\sum_{i=1}^{N} \mu_i \left( \frac{\vect{r} - \vect{r}_i}{\norm{\vect{r} - \vect{r}_i}^3} + \frac{\vect{r}_i}{\norm{\vect{r}_i}^3} \right)
    \label{eq:nbody}
\end{equation}

where:
\begin{itemize}
    \item $\vect{r}$ is the spacecraft position in SSB frame
    \item $\vect{r}_i$ is the position of body $i$ in SSB frame
    \item $\mu_i = GM_i$ is the gravitational parameter of body $i$
\end{itemize}

The second term in \cref{eq:nbody} accounts for the indirect acceleration due to the motion of the SSB relative to each body.

\subsection{Two-Body Problem}

For simplified analysis relative to a central body, the equation of motion reduces to:

\begin{equation}
    \ddot{\vect{r}}_\text{rel} = -\frac{\mu}{\norm{\vect{r}_\text{rel}}^3} \vect{r}_\text{rel}
    \label{eq:twobody}
\end{equation}

\subsection{Gravitational Parameters}

The simulation uses the following gravitational parameters (in $\si{AU^3.day^{-2}}$):

\begin{table}[h]
\centering
\begin{tabular}{lcc}
\toprule
\textbf{Body} & $\mu$ [$\si{AU^3.day^{-2}}$] & $\mu$ [$\si{km^3.s^{-2}}$] \\
\midrule
Sun & $2.959122 \times 10^{-4}$ & $1.327124 \times 10^{11}$ \\
Earth & $8.887692 \times 10^{-10}$ & $3.986004 \times 10^{5}$ \\
Moon & $1.093189 \times 10^{-11}$ & $4.902800 \times 10^{3}$ \\
Mars & $9.549548 \times 10^{-11}$ & $4.282837 \times 10^{4}$ \\
Jupiter & $2.824761 \times 10^{-7}$ & $1.266865 \times 10^{8}$ \\
\bottomrule
\end{tabular}
\caption{Gravitational parameters for major bodies.}
\label{tab:mu}
\end{table}

%==============================================================================
\section{Orbital Mechanics}
%==============================================================================

\subsection{Keplerian Elements}

The classical orbital elements $(\texttt{a}, \texttt{e}, \texttt{i}, \Omega, \omega, \nu)$ are defined as:

\begin{itemize}
    \item $a$ -- Semi-major axis
    \item $e$ -- Eccentricity
    \item $i$ -- Inclination
    \item $\Omega$ -- Right ascension of ascending node (RAAN)
    \item $\omega$ -- Argument of periapsis
    \item $\nu$ -- True anomaly
\end{itemize}

\subsection{Vis-Viva Equation}

The orbital velocity at radius $r$ in an orbit with semi-major axis $a$:

\begin{equation}
    v = \sqrt{\mu \left( \frac{2}{r} - \frac{1}{a} \right)}
    \label{eq:visviva}
\end{equation}

For circular orbits ($a = r$):
\begin{equation}
    v_\text{circ} = \sqrt{\frac{\mu}{r}}
\end{equation}

For escape velocity ($a \to \infty$):
\begin{equation}
    v_\text{esc} = \sqrt{\frac{2\mu}{r}}
\end{equation}

\subsection{Hohmann Transfer}

For a coplanar transfer between circular orbits of radii $r_1$ and $r_2$, the transfer orbit has semi-major axis:

\begin{equation}
    a_\text{trans} = \frac{r_1 + r_2}{2}
\end{equation}

The velocity changes at periapsis (burn 1) and apoapsis (burn 2) are:

\begin{align}
    \Delta v_1 &= v_\text{trans,peri} - v_{\text{circ},1} = \sqrt{\frac{\mu}{r_1}} \left( \sqrt{\frac{2r_2}{r_1 + r_2}} - 1 \right) \\
    \Delta v_2 &= v_{\text{circ},2} - v_\text{trans,apo} = \sqrt{\frac{\mu}{r_2}} \left( 1 - \sqrt{\frac{2r_1}{r_1 + r_2}} \right)
\end{align}

The total $\Delta v$ for a Hohmann transfer:
\begin{equation}
    \Delta v_\text{total} = \abs{\Delta v_1} + \abs{\Delta v_2}
\end{equation}

\subsection{Time of Flight}

The transfer time for a Hohmann maneuver (half the transfer orbit period):

\begin{equation}
    t_\text{transfer} = \pi \sqrt{\frac{a_\text{trans}^3}{\mu}}
\end{equation}

\subsection{Orbital Period}

The period of an elliptical orbit:

\begin{equation}
    T = 2\pi \sqrt{\frac{a^3}{\mu}}
\end{equation}

%==============================================================================
\section{Attitude Dynamics}
%==============================================================================

\subsection{Quaternion Representation}

Attitude is represented using unit quaternions $\quat{q} = [q_0, q_1, q_2, q_3]^T$ where:
\begin{equation}
    \quat{q} = q_0 + q_1 \vect{i} + q_2 \vect{j} + q_3 \vect{k}
\end{equation}

with the constraint:
\begin{equation}
    \norm{\quat{q}}^2 = q_0^2 + q_1^2 + q_2^2 + q_3^2 = 1
\end{equation}

\subsection{Quaternion Operations}

\subsubsection{Quaternion Multiplication}

For quaternions $\quat{p}$ and $\quat{q}$:

\begin{equation}
    \quat{p} \otimes \quat{q} =
    \begin{bmatrix}
        p_0 q_0 - p_1 q_1 - p_2 q_2 - p_3 q_3 \\
        p_0 q_1 + p_1 q_0 + p_2 q_3 - p_3 q_2 \\
        p_0 q_2 - p_1 q_3 + p_2 q_0 + p_3 q_1 \\
        p_0 q_3 + p_1 q_2 - p_2 q_1 + p_3 q_0
    \end{bmatrix}
\end{equation}

\subsubsection{Quaternion Conjugate}

\begin{equation}
    \quat{q}^* = [q_0, -q_1, -q_2, -q_3]^T
\end{equation}

\subsubsection{Direction Cosine Matrix}

The rotation matrix from body frame to inertial frame:

\begin{equation}
    \mat{R} = \begin{bmatrix}
        1 - 2(q_2^2 + q_3^2) & 2(q_1 q_2 - q_0 q_3) & 2(q_1 q_3 + q_0 q_2) \\
        2(q_1 q_2 + q_0 q_3) & 1 - 2(q_1^2 + q_3^2) & 2(q_2 q_3 - q_0 q_1) \\
        2(q_1 q_3 - q_0 q_2) & 2(q_2 q_3 + q_0 q_1) & 1 - 2(q_1^2 + q_2^2)
    \end{bmatrix}
\end{equation}

\subsection{Quaternion Kinematics}

The quaternion derivative given angular velocity $\vect{\omega}$ (in body frame):

\begin{equation}
    \dot{\quat{q}} = \frac{1}{2} \mat{\Omega}(\vect{\omega}) \quat{q}
    \label{eq:qdot}
\end{equation}

where:
\begin{equation}
    \mat{\Omega}(\vect{\omega}) = \begin{bmatrix}
        0 & -\omega_x & -\omega_y & -\omega_z \\
        \omega_x & 0 & \omega_z & -\omega_y \\
        \omega_y & -\omega_z & 0 & \omega_x \\
        \omega_z & \omega_y & -\omega_x & 0
    \end{bmatrix}
\end{equation}

\subsection{Euler's Equations}

The rotational dynamics of a rigid body:

\begin{equation}
    \mat{I} \dot{\vect{\omega}} + \vect{\omega} \times (\mat{I} \vect{\omega}) = \vect{\tau}_\text{ext}
    \label{eq:euler}
\end{equation}

where:
\begin{itemize}
    \item $\mat{I}$ is the inertia tensor
    \item $\vect{\omega}$ is the angular velocity in body frame
    \item $\vect{\tau}_\text{ext}$ is the external torque
\end{itemize}

For a diagonal inertia tensor $\mat{I} = \text{diag}(I_x, I_y, I_z)$:

\begin{align}
    I_x \dot{\omega}_x &= (I_y - I_z) \omega_y \omega_z + \tau_x \\
    I_y \dot{\omega}_y &= (I_z - I_x) \omega_z \omega_x + \tau_y \\
    I_z \dot{\omega}_z &= (I_x - I_y) \omega_x \omega_y + \tau_z
\end{align}

%==============================================================================
\section{Attitude Control}
%==============================================================================

\subsection{PD Controller}

The proportional-derivative controller generates commanded torque:

\begin{equation}
    \vect{\tau}_\text{cmd} = -K_p \vect{q}_\text{err} - K_d \vect{\omega}
    \label{eq:pd}
\end{equation}

where:
\begin{itemize}
    \item $K_p$ is the proportional gain
    \item $K_d$ is the derivative gain
    \item $\vect{q}_\text{err} = [q_1, q_2, q_3]^T$ is the vector part of the quaternion error
    \item $\vect{\omega}$ is the angular velocity
\end{itemize}

The quaternion error is computed as:
\begin{equation}
    \quat{q}_\text{err} = \quat{q}_\text{des}^* \otimes \quat{q}
\end{equation}

\subsection{Reaction Wheel Control}

The simulation uses a tetrahedral 4-wheel configuration with wheel axes:

\begin{equation}
    \mat{A}_\text{wheel} = \begin{bmatrix}
        1 & -1 & -1 & 1 \\
        1 & 1 & -1 & -1 \\
        1 & 1 & 1 & 1
    \end{bmatrix} \cdot \frac{1}{\sqrt{3}}
\end{equation}

Wheel torque allocation uses the pseudoinverse:

\begin{equation}
    \vect{\tau}_\text{wheel} = \mat{A}_\text{wheel}^+ \vect{\tau}_\text{cmd}
\end{equation}

where $\mat{A}^+ = \mat{A}^T (\mat{A} \mat{A}^T)^{-1}$ is the Moore-Penrose pseudoinverse.

\subsubsection{Momentum Management}

Each wheel has angular momentum:
\begin{equation}
    h_i = I_\text{wheel} \Omega_i
\end{equation}

Total spacecraft angular momentum from wheels:
\begin{equation}
    \vect{H}_\text{wheel} = \mat{A}_\text{wheel} \vect{h}
\end{equation}

Saturation occurs when $|h_i| > h_\text{max}$.

\subsection{RCS Thruster Control}

For situations requiring high torque or momentum desaturation, the simulation employs a 16-thruster RCS configuration.

The torque from thruster $i$ is:
\begin{equation}
    \vect{\tau}_i = \vect{r}_i \times (F_i \hat{\vect{n}}_i)
\end{equation}

where:
\begin{itemize}
    \item $\vect{r}_i$ is the thruster position relative to center of mass
    \item $F_i$ is the thrust magnitude
    \item $\hat{\vect{n}}_i$ is the thrust direction
\end{itemize}

The torque allocation matrix $\mat{B}$ maps thruster firings to body torques:
\begin{equation}
    \vect{\tau} = \mat{B} \vect{u}
\end{equation}

where $\vect{u} \in [0,1]^{16}$ is the thruster duty cycle vector.

\subsubsection{Mass Flow Rate}

RCS propellant consumption:
\begin{equation}
    \dot{m}_\text{RCS} = \frac{\sum_i F_i u_i}{I_\text{sp,RCS} \cdot g_0}
\end{equation}

%==============================================================================
\section{Propulsion}
%==============================================================================

\subsection{Thrust Model}

Thrust force from a rocket engine:
\begin{equation}
    F = \dot{m} v_e = \dot{m} I_\text{sp} g_0
\end{equation}

where:
\begin{itemize}
    \item $\dot{m}$ is the mass flow rate
    \item $v_e$ is the exhaust velocity
    \item $I_\text{sp}$ is the specific impulse
    \item $g_0 = \SI{9.80665}{m.s^{-2}}$ is standard gravity
\end{itemize}

\subsection{Tsiolkovsky Rocket Equation}

The ideal velocity change for a given propellant mass:
\begin{equation}
    \Delta v = v_e \ln \frac{m_0}{m_f} = I_\text{sp} g_0 \ln \frac{m_0}{m_f}
    \label{eq:rocket}
\end{equation}

Rearranging for propellant mass:
\begin{equation}
    m_\text{prop} = m_0 \left( 1 - e^{-\Delta v / v_e} \right)
\end{equation}

\subsection{Burn Duration}

For a constant-thrust maneuver:
\begin{equation}
    t_\text{burn} = \frac{m_\text{prop}}{\dot{m}} = \frac{m_0 (1 - e^{-\Delta v / v_e})}{\dot{m}}
\end{equation}

\subsection{Equations of Motion with Thrust}

The full equations of motion including thrust:

\begin{align}
    \dot{\vect{r}} &= \vect{v} \\
    \dot{\vect{v}} &= \vect{a}_\text{grav} + \frac{F}{m} \hat{\vect{u}} \\
    \dot{m} &= -\frac{F}{I_\text{sp} g_0}
\end{align}

where $\hat{\vect{u}}$ is the thrust direction unit vector.

%==============================================================================
\section{Guidance Algorithms}
%==============================================================================

\subsection{ZEM/ZEV Guidance}

The Zero-Effort-Miss/Zero-Effort-Velocity (ZEM/ZEV) guidance algorithm provides closed-loop thrust direction for terminal guidance.

\subsubsection{Zero-Effort-Miss}

The position error that would result from zero thrust:
\begin{equation}
    \vect{\text{ZEM}} = \vect{r}_\text{pred}(t_f) - \vect{r}_\text{target}
\end{equation}

\subsubsection{Zero-Effort-Velocity}

The velocity error at target:
\begin{equation}
    \vect{\text{ZEV}} = \vect{v}_\text{pred}(t_f) - \vect{v}_\text{target}
\end{equation}

\subsubsection{Optimal Steering}

The optimal thrust direction combines ZEM and ZEV:
\begin{equation}
    \vect{\Delta v}_\text{req} = \frac{\vect{\text{ZEM}}}{t_\text{go}} + \vect{\text{ZEV}}
\end{equation}

where $t_\text{go} = t_f - t$ is time-to-go.

The thrust direction:
\begin{equation}
    \hat{\vect{u}} = \frac{\vect{\Delta v}_\text{req}}{\norm{\vect{\Delta v}_\text{req}}}
\end{equation}

\subsubsection{Velocity-to-Go}

The remaining velocity change:
\begin{equation}
    v_\text{go} = \norm{\vect{\Delta v}_\text{req}}
\end{equation}

Engine cutoff occurs when $v_\text{go} < v_\text{threshold}$.

\subsection{Finite Burn Implementation}

The ZEM/ZEV algorithm iteratively refines the solution:

\begin{quote}
\textbf{Algorithm: ZEM/ZEV Guidance Update}
\begin{enumerate}
    \item Initialize $\hat{\vect{u}} \gets \hat{\vect{v}}$ (velocity direction)
    \item For $k = 1$ to $N_\text{iter}$:
    \begin{enumerate}
        \item Propagate with constant $\hat{\vect{u}}$ to $t_f$
        \item Compute $\vect{\text{ZEM}} = \vect{r}_\text{pred} - \vect{r}_\text{target}$
        \item Compute $\vect{\text{ZEV}} = \vect{v}_\text{pred} - \vect{v}_\text{target}$
        \item $\vect{\Delta v}_\text{req} = \vect{\text{ZEM}}/t_\text{go} + \vect{\text{ZEV}}$
        \item Update: $\hat{\vect{u}} \gets \alpha \hat{\vect{u}}_\text{new} + (1-\alpha) \hat{\vect{u}}$ (damping)
    \end{enumerate}
    \item Return $v_\text{go}$, $\hat{\vect{u}}$
\end{enumerate}
\end{quote}

The damping factor $\alpha \approx 0.6$ prevents oscillation.

%==============================================================================
\section{Numerical Integration}
%==============================================================================

\subsection{State Vector}

The full state vector for propagation:

\begin{equation}
    \vect{y} = [\vect{r}, \vect{v}, m, \quat{q}, \vect{\omega}, \vect{h}_\text{wheel}]^T
\end{equation}

Dimensions:
\begin{itemize}
    \item Position $\vect{r}$: 3 components (AU)
    \item Velocity $\vect{v}$: 3 components (AU/day)
    \item Mass $m$: 1 component (kg)
    \item Quaternion $\quat{q}$: 4 components (normalized)
    \item Angular velocity $\vect{\omega}$: 3 components (rad/day)
    \item Wheel momentum $\vect{h}_\text{wheel}$: 4 components (N$\cdot$m$\cdot$s)
\end{itemize}

Total: 18 state variables.

\subsection{Integration Method}

The simulation uses the DOP853 integrator (8th-order Dormand-Prince with dense output) with:

\begin{itemize}
    \item Relative tolerance: $10^{-9}$
    \item Absolute tolerance: $10^{-9}$
    \item Maximum step size: $\SI{1}{\minute}$ (for attitude accuracy)
\end{itemize}

\subsection{Quaternion Normalization}

After each integration step, the quaternion is renormalized:

\begin{equation}
    \quat{q} \gets \frac{\quat{q}}{\norm{\quat{q}}}
\end{equation}

This maintains the unit constraint to machine precision ($\sim 10^{-16}$).

%==============================================================================
\section{Validation Results}
%==============================================================================

\subsection{Orbital Mechanics Accuracy}

\begin{table}[h]
\centering
\begin{tabular}{lcc}
\toprule
\textbf{Metric} & \textbf{Target} & \textbf{Achieved} \\
\midrule
Position drift (1 orbit) & $< \SI{1}{km}$ & $< \SI{0.002}{km}$ \\
Quaternion norm error & $< 10^{-6}$ & $\sim 10^{-16}$ \\
Angular rate stability & $< \SI{0.1}{\degree/s}$ & $\sim \SI{0}{\degree/s}$ \\
Hohmann $\Delta v$ accuracy & $< 1\%$ & $< 0.1\%$ \\
\bottomrule
\end{tabular}
\caption{Numerical accuracy validation results.}
\end{table}

\subsection{Attitude Control Performance}

\begin{table}[h]
\centering
\begin{tabular}{lcc}
\toprule
\textbf{Test} & \textbf{Target} & \textbf{Achieved} \\
\midrule
Mean pointing error & $< \SI{5}{\degree}$ & $\SI{0.11}{\degree}$ \\
Max pointing error & $< \SI{10}{\degree}$ & $\SI{0.77}{\degree}$ \\
Settling time & $< \SI{60}{s}$ & $< \SI{10}{s}$ \\
\bottomrule
\end{tabular}
\caption{Attitude control validation results.}
\end{table}

%==============================================================================
\section{Conclusion}
%==============================================================================

This document has presented the mathematical foundations of the Space State astrodynamics simulation. The implementation achieves high numerical accuracy through:

\begin{enumerate}
    \item High-fidelity n-body gravitational dynamics
    \item Quaternion-based attitude representation avoiding gimbal lock
    \item Adaptive integration with tight tolerances
    \item Validated guidance algorithms for finite burns
\end{enumerate}

All physics models have been verified through comprehensive unit testing (616 tests) and mission validation scenarios.

%==============================================================================
\appendix
\section{Physical Constants}
%==============================================================================

\begin{table}[h]
\centering
\begin{tabular}{llc}
\toprule
\textbf{Constant} & \textbf{Symbol} & \textbf{Value} \\
\midrule
Gravitational constant & $G$ & $\SI{6.67430e-11}{m^3.kg^{-1}.s^{-2}}$ \\
Speed of light & $c$ & $\SI{299792458}{m.s^{-1}}$ \\
Astronomical unit & AU & $\SI{149597870.7}{km}$ \\
Solar mass & $M_\odot$ & $\SI{1.98892e30}{kg}$ \\
Earth mass & $M_\oplus$ & $\SI{5.9722e24}{kg}$ \\
Earth radius & $R_\oplus$ & $\SI{6371}{km}$ \\
Moon mass & $M_\text{Moon}$ & $\SI{7.342e22}{kg}$ \\
Moon radius & $R_\text{Moon}$ & $\SI{1737.4}{km}$ \\
Standard gravity & $g_0$ & $\SI{9.80665}{m.s^{-2}}$ \\
\bottomrule
\end{tabular}
\caption{Physical constants.}
\end{table}

%==============================================================================
\section{Spacecraft Parameters}
%==============================================================================

Reference spacecraft (Starship configuration):

\begin{table}[h]
\centering
\begin{tabular}{lc}
\toprule
\textbf{Parameter} & \textbf{Value} \\
\midrule
Dry mass & $\SI{120000}{kg}$ \\
Propellant mass & $\SI{1500000}{kg}$ \\
RCS propellant & $\SI{2000}{kg}$ \\
Main engine thrust & $\SI{2200}{kN}$ \\
Main engine $I_\text{sp}$ & $\SI{380}{s}$ \\
Inertia (x, y, z) & $\SI{500000}{}, \SI{500000}{}, \SI{100000}{kg.m^2}$ \\
Wheel momentum capacity & $\SI{100}{N.m.s}$ \\
\bottomrule
\end{tabular}
\caption{Reference spacecraft parameters.}
\end{table}

\end{document}
